\documentclass{beamer}
\usetheme{Singapore}

\usepackage[T2A]{fontenc}
\usepackage[utf8]{inputenc}
\usepackage[ukrainian]{babel}

\usepackage{copyrightbox}
\usepackage{tikz}
\usepackage{pgf-umlcd}
\usepackage{pgfplots}

\title{Моделювання процесів термопружності у композитах методами граничних і приграничних елементів}
\author{Андрій Петльований}
\institute{Львівський національний університет ім. І. Франка}
\date{\today}

\begin{document}

\begin{frame}
    \titlepage
\end{frame}

\begin{frame}
    \frametitle{Порядок денний}
    \tableofcontents
\end{frame}

\section{Вступ}

\section{Задача теплопровідності}
% \subsection{sub a}

\begin{frame}
    \frametitle{Постановка задачі}
    \begin{equation}
        \label{eqn:theory_thermo_pennes_eqn}
        \nabla \cdot \lambda_\theta \nabla \theta(x) + \omega \rho_b C_b (\theta_c - \theta(x)) + q_m  = 0
    \end{equation}
\end{frame}

\begin{frame}
    \frametitle{Область}
    \begin{figure}[ht!]
        \centering
        \begin{tikzpicture}[scale=0.75, transform shape]
            \draw (1,1) rectangle (11,3);
            \draw (6,2.5) circle (0.5);
      
            \draw[thick,->] (0,0) -- (12,0) node[anchor=north west] {x};
            \draw[thick,->] (0,0) -- (0,6) node[anchor=south east] {y};        
        
            \foreach \x in {0,...,4}
                \pgfmathtruncatemacro{\label}{\x * 10}
                \draw ( \x * 2.5 cm,1pt) -- ( \x * 2.5 cm,-1pt) node[anchor=north] {\label};
            \foreach \y in {0,...,2}
                \pgfmathtruncatemacro{\label}{\y * 10}
                \draw (1pt,\y * 2.5 cm) -- (-1pt,\y * 2.5 cm) node[anchor=east] {\label};
            
            \node at (8,2) {$\Omega$};
            \node at (6,2.5) {$H$};
            \node at (4,0.5) {$\Gamma$};
      
        \end{tikzpicture}
      \end{figure}
\end{frame}

\begin{frame}
    \frametitle{Коефіцієнт теплопровідності}
    \begin{figure}[ht!]
        \centering
        \begin{tikzpicture}[domain=-5:5, scale=0.75, transform shape]
            \draw[->] (-7,0) -- (7,0) node[right] {$x_1$};
            \draw[->] (0,-1) -- (0,5) node[above] {$5\lambda_\theta(x)$};
        
            \draw[color=blue] plot (\x, {5 * (0.305 * cos(1.57 / 25 * (\x * \x r)) + 0.19)});
            \draw[color=blue] (-6,0.95) -- (-5,0.95);
            \draw[color=blue] (5,0.95) -- (6,0.95);
        \end{tikzpicture}
    \end{figure}
\end{frame}

\begin{frame}
    \frametitle{Адитивне розщеплення оператора}
    \begin{equation}
        \label{eqn:thermo_pennes_separation_eqn}
        \lambda_{\theta c} \nabla^2 \theta(x) + \omega \rho_b C_b (\theta_c - \theta(x)) = 
          - \nabla \lambda_{\theta H} (x) \cdot \nabla \theta(x) \cdot \chi_H
      \end{equation}
\end{frame}

\begin{frame}
    \frametitle{МГЕ -- інтегральне подання розв'язку}
    \begin{equation}
        \label{eqn:thermo_bem_direct_pennes}
        \begin{split}
        c(\xi)\theta(\xi) + \int \limits_{\Gamma}\theta(x)\hat{q}(x,\xi)d\Gamma(x) + \\
            + \int \limits_{H} \nabla{\lambda_{\theta H}(x)}\nabla{\theta(x)} \hat{\theta}(x,\xi) d H(x) 
            = \int \limits_{\Gamma}q(x)\hat{\theta}(x,\xi)d\Gamma(x)
        \end{split}
      \end{equation}
\end{frame}

\begin{frame}
    \frametitle{МГЕ -- зведення за формулою Гріна}
    \begin{equation}
        \label{eqn:thermo_bem_direct_pennes_simplified}
        \begin{split}
          c(\xi)\theta(\xi) + \int \limits_{\Gamma}\theta(x)\hat{q}(x,\xi)d\Gamma(x) = \\
            = \int \limits_{H}\theta(x) [ \nabla^2 \lambda_{\theta H}(x) \hat{\theta}(x,\xi) + 
                \nabla \lambda_{\theta H}(x) \nabla \hat{\theta}(x,\xi) ] dH(x) + \\
            + \int \limits_{\Gamma}q(x)\hat{\theta}(x,\xi)d\Gamma(x)
        \end{split}
      \end{equation}
\end{frame}

\begin{frame}
    \frametitle{Висновки}
    Lorem ipsum dolor sit amet, consectetur adipisicing elit, sed do eiusmod tempor incididunt ut labore et dolore magna
    aliqua.
\end{frame}


\section{Задача пружності}

\begin{frame}
    \frametitle{Title}
    Lorem ipsum dolor sit amet, consectetur adipisicing elit, sed do eiusmod tempor incididunt ut labore et dolore magna
    aliqua.
\end{frame}

\begin{frame}
    \frametitle{Область}
    \begin{figure}[ht!]
        \centering
        \begin{tikzpicture}[scale=0.75, transform shape]
            \draw (5.7,3) rectangle (6.3,7);
            \draw (1,1) rectangle (11,3);
            \draw (6,2.5) circle (0.5);
            \draw[thick, ->] (5.3,6) -- (5.3,4);
      
            \draw[thick,->] (0,0) -- (12,0) node[anchor=north west] {x};
            \draw[thick,->] (0,0) -- (0,9) node[anchor=south east] {y};        
        
            \foreach \x in {0,...,4}
                \pgfmathtruncatemacro{\label}{\x * 10}
                \draw ( \x * 2.5 cm,1pt) -- ( \x * 2.5 cm,-1pt) node[anchor=north] {\label};
            \foreach \y in {0,...,3}
                \pgfmathtruncatemacro{\label}{\y * 10}
                \draw (1pt,\y * 2.5 cm) -- (-1pt,\y * 2.5 cm) node[anchor=east] {\label};
            
            \node at (6,3.5) {$\Gamma_i$};
      
            \node at (8,2) {$\Omega$};
            \node at (6,5) {$\Omega_i$};
            \node at (6,2.5) {$H$};
            \node at (4,0.5) {$\Gamma_b$};
      
        \end{tikzpicture}
      \end{figure}
\end{frame}

\section{Задача термопружності}

\begin{frame}
    \frametitle{Title}
    Lorem ipsum dolor sit amet, consectetur adipisicing elit, sed do eiusmod tempor incididunt ut labore et dolore magna
    aliqua.
\end{frame}

\section{Підсумки}

\begin{frame}
    \frametitle{Title}
    Lorem ipsum dolor sit amet, consectetur adipisicing elit, sed do eiusmod tempor incididunt ut labore et dolore magna
    aliqua.
\end{frame}


\end{document}

